%%=============================================================================
%% Woord vooraf
%%=============================================================================

\chapter*{\IfLanguageName{dutch}{Woord vooraf}{Preface}}%
\label{ch:voorwoord}
TODO!
%% Richtlengte: een halve bladzijde. Dit is het enige deel van je rapport waar
%% je in de ik-vorm schrijft.
%% Schrijf hier kort:
%% - waarom dit onderwerp jou interesseerde;
%% - wat je ervan verwachtte toen je eraan begon;
%% - wie je wil bedanken (je mentor, je team, je begeleider).
%% Let op: dit is geen samenvatting van je project. Wie enkel dit leest, hoeft
%% nog niet te weten wat je gebouwd hebt.
Mijn graduaatsproject vormde het sluitstuk van mijn opleiding Graduaat in het Programmeren aan HOGENT en werd uitgevoerd tijdens mijn stage bij Vesalius.ai. Tijdens deze stage kreeg ik de mogelijkheid om mee te werken aan een professioneel softwareproduct en kennis te maken met een ontwikkelomgeving die aanzienlijk uitgebreider is dan de projecten waarmee ik tijdens de opleiding had gewerkt.
Ik wil in de eerste plaats Sander Goossens, CEO van Vesalius.ai, bedanken voor de mogelijkheid om mijn stage binnen het bedrijf uit te voeren en voor het vertrouwen dat mij werd gegeven om aan een concreet onderdeel van het platform mee te werken.
Daarnaast wil ik Bart D'Hoore, CSM en Head of People, bedanken voor de introductie, onboarding en begeleiding tijdens mijn start binnen het bedrijf.
Ook wil ik het volledige Vesalius.ai-team bedanken voor de aangename samenwerking, de vlotte communicatie en de bereidheid om vragen te beantwoorden en ondersteuning te bieden tijdens mijn stage. In het bijzonder wil ik Quinten Schietecatte, Navaron Bracke, Ruben Redant en Axel Jonckheere bedanken voor hun hulp en technische ondersteuning tijdens het project.
Tot slot heeft deze stage mij de mogelijkheid gegeven om de kennis uit mijn opleiding toe te passen binnen een echte softwareomgeving en tegelijk te leren omgaan met een bestaande codebase, professionele ontwikkelprocessen en de verwachtingen van een developmentteam.
